\documentclass[tikz,border=6pt]{standalone}
\usepackage{kotex}
\usetikzlibrary{arrows.meta,positioning}
\begin{document}
\begin{tikzpicture}[
  font=\sffamily,
  state/.style={draw=blue!65!black, fill=blue!5, rounded corners=3pt,
    minimum width=69mm, minimum height=12mm, align=center},
  command/.style={font=\ttfamily\small, fill=white, inner sep=2pt},
  flow/.style={-{Stealth[length=2.5mm]}, thick, draw=black!75},
  node distance=11mm
]
\node[state] (work) {\textbf{작업 폴더}\\[-1pt]\small 편집한 파일};
\node[state, below=of work] (stage) {\textbf{스테이징 영역}\\[-1pt]\small 다음 커밋에 담을 파일 내용};
\node[state, below=of stage] (local) {\textbf{로컬 저장소}\\[-1pt]\small 내 컴퓨터에 기록된 커밋};
\node[state, draw=teal!65!black, fill=teal!6, below=of local] (remote)
  {\textbf{GitHub 원격 저장소}\\[-1pt]\small 온라인에 올라간 커밋};
\draw[flow] (work) -- node[command] {git add} (stage);
\draw[flow] (stage) -- node[command] {git commit} (local);
\draw[flow] (local) -- node[command] {git push} (remote);
\end{tikzpicture}
\end{document}
